% --- Archon physics formalization source begin ---
% archon:physics
% archon:covers PhyXMiniProblems/problem_phyx_mini_0251.lean
% archon:source-report reports/phyx_mini/problem_phyx_mini_0251.source.json

\chapter{Physics problem phyx\_mini\_0251}
\label{ch:PhyXMiniProblems_problem_phyx_mini_0251}

\paragraph{Problem source.}
\#\# Physical scenario

The vertical axis scale is set by \$x\_s = 6.0\textbackslash{} cm\$.

\#\# Question

What is the phase constant (positve) for the harmonic oscillator with the position function \$x(t)\$ given in figure if the position function has the form \$x = x\_m \textbackslash{}cos(\textbackslash{}omega t + \textbackslash{}phi)\$?

\#\# Answer choices

- A: A: 1.15 rad

- B: B: 1.47 rad

- C: C: 1.91 rad

- D: D: 2.23 rad

\#\# Figure caption (auxiliary; use the image as primary evidence)

The image is a graph showing a sinusoidal wave. Here's a detailed description of its components:

- The horizontal axis is labeled \textbackslash{}( t \textbackslash{}), likely representing time.
- The vertical axis is labeled \textbackslash{}( x \textbackslash{}) with units in centimeters (cm).
- The vertical axis has specific points marked as \textbackslash{}( x\_s \textbackslash{}) and \textbackslash{}(-x\_s\textbackslash{}).
- The wave moves from top left to bottom right, characteristic of a sine or cosine wave.
- Horizontal grid lines are spaced along the graph, potentially for reference or measuring purposes.
- The graph includes thick lines emphasizing the wave's curve and the axes.

\#\# Dataset metadata

Wave Properties; Physical Model Grounding Reasoning, Spatial Relation Reasoning

\paragraph{Recorded answer/context.}
C: C: 1.91 rad

\paragraph{Figure/image path.}
/root/proposal\_for\_physic/science-mango/phyx\_mini\_run/phyx\_data/test\_image/251.png

\paragraph{Formalization target.}
create a compiling Lean file with sorry bodies at `PhyXMiniProblems/problem\_phyx\_mini\_0251.lean`. The Lean declarations must preserve the physical quantities, dimensions or dimensional roles, figure labels, governing-law hypotheses, and final relation expressed by this problem.
Use Mathlib/Physlib names found through LeanExplore where available. If a physics API is missing, introduce faithful local abstractions rather than scalar placeholder aliases.

\begin{theorem}[Physics formalization target]
\label{thm:physics:phyx_mini_0251:target}
\lean{PhyXMiniProblems.ProblemPhyXMini0251.problem_phyx_mini_0251}
\uses{def:physics:phyx-mini-0251:phyxminiproblems-problemphyxmini0251-harmonicoscillatorgraphsetup, def:physics:phyx-mini-0251:phyxminiproblems-problemphyxmini0251-matchesoscillatorfigure, def:physics:phyx-mini-0251:phyxminiproblems-problemphyxmini0251-hasphysicalharmonicparameters, def:physics:phyx-mini-0251:phyxminiproblems-problemphyxmini0251-satisfiescosineharmonicmotion, def:physics:phyx-mini-0251:phyxminiproblems-problemphyxmini0251-answerchoice, def:physics:phyx-mini-0251:phyxminiproblems-problemphyxmini0251-answerchoiceradians, def:physics:phyx-mini-0251:phyxminiproblems-problemphyxmini0251-recordedanswerchoice, def:physics:phyx-mini-0251:phyxminiproblems-problemphyxmini0251-matchesanswertonearesthundredth}
The assigned autoformalize agent should translate this physics problem into Lean declarations in the covered file, with theorem and lemma proof bodies written as `by sorry`.
\end{theorem}
\begin{proof}
This is an autoformalization task, not a proof task. Produce statements that can later be proved by the physics prover without weakening the physical model.
\end{proof}

% --- Archon declaration topology begin ---
\section{Declaration topology}
\begin{definition}[Length Quantity]
  \label{def:physics:phyx-mini-0251:phyxminiproblems-problemphyxmini0251-lengthquantity}
  \lean{PhyXMiniProblems.ProblemPhyXMini0251.LengthQuantity}
  A signed one-dimensional displacement or oscillation amplitude
\end{definition}
\begin{proof}
  This is the declared model definition of Length Quantity.
\end{proof}
\begin{definition}[Time Quantity]
  \label{def:physics:phyx-mini-0251:phyxminiproblems-problemphyxmini0251-timequantity}
  \lean{PhyXMiniProblems.ProblemPhyXMini0251.TimeQuantity}
  A physical time coordinate
\end{definition}
\begin{proof}
  This is the declared model definition of Time Quantity.
\end{proof}
\begin{definition}[Angular Frequency Quantity]
  \label{def:physics:phyx-mini-0251:phyxminiproblems-problemphyxmini0251-angularfrequencyquantity}
  \lean{PhyXMiniProblems.ProblemPhyXMini0251.AngularFrequencyQuantity}
  Angular frequency; radians are dimensionless, so this has inverse-time dimension
\end{definition}
\begin{proof}
  This is the declared model definition of Angular Frequency Quantity.
\end{proof}
\begin{definition}[Velocity Quantity]
  \label{def:physics:phyx-mini-0251:phyxminiproblems-problemphyxmini0251-velocityquantity}
  \lean{PhyXMiniProblems.ProblemPhyXMini0251.VelocityQuantity}
  A signed one-dimensional velocity
\end{definition}
\begin{proof}
  This is the declared model definition of Velocity Quantity.
\end{proof}
\begin{definition}[length Readout]
  \label{def:physics:phyx-mini-0251:phyxminiproblems-problemphyxmini0251-lengthreadout}
  \lean{PhyXMiniProblems.ProblemPhyXMini0251.lengthReadout}
  \uses{def:physics:phyx-mini-0251:phyxminiproblems-problemphyxmini0251-lengthquantity}
  Read a physical length in the selected length unit
\end{definition}
\begin{proof}
  This is the declared model definition of length Readout.
\end{proof}
\begin{definition}[time Readout]
  \label{def:physics:phyx-mini-0251:phyxminiproblems-problemphyxmini0251-timereadout}
  \lean{PhyXMiniProblems.ProblemPhyXMini0251.timeReadout}
  \uses{def:physics:phyx-mini-0251:phyxminiproblems-problemphyxmini0251-timequantity}
  Read a physical time in the selected time unit
\end{definition}
\begin{proof}
  This is the declared model definition of time Readout.
\end{proof}
\begin{definition}[angular Frequency Readout]
  \label{def:physics:phyx-mini-0251:phyxminiproblems-problemphyxmini0251-angularfrequencyreadout}
  \lean{PhyXMiniProblems.ProblemPhyXMini0251.angularFrequencyReadout}
  \uses{def:physics:phyx-mini-0251:phyxminiproblems-problemphyxmini0251-angularfrequencyquantity}
  Read angular frequency in radians per selected time unit
\end{definition}
\begin{proof}
  This is the declared model definition of angular Frequency Readout.
\end{proof}
\begin{definition}[velocity Readout]
  \label{def:physics:phyx-mini-0251:phyxminiproblems-problemphyxmini0251-velocityreadout}
  \lean{PhyXMiniProblems.ProblemPhyXMini0251.velocityReadout}
  \uses{def:physics:phyx-mini-0251:phyxminiproblems-problemphyxmini0251-velocityquantity}
  Read velocity in the coherent selected length/time unit system
\end{definition}
\begin{proof}
  This is the declared model definition of velocity Readout.
\end{proof}
\begin{definition}[Graph Axis]
  \label{def:physics:phyx-mini-0251:phyxminiproblems-problemphyxmini0251-graphaxis}
  \lean{PhyXMiniProblems.ProblemPhyXMini0251.GraphAxis}
  The two coordinate axes shown in the graph
\end{definition}
\begin{proof}
  This is the declared model definition of Graph Axis.
\end{proof}
\begin{definition}[Axis Quantity]
  \label{def:physics:phyx-mini-0251:phyxminiproblems-problemphyxmini0251-axisquantity}
  \lean{PhyXMiniProblems.ProblemPhyXMini0251.AxisQuantity}
  The physical quantity printed on each graph axis
\end{definition}
\begin{proof}
  This is the declared model definition of Axis Quantity.
\end{proof}
\begin{definition}[Displacement Time Graph]
  \label{def:physics:phyx-mini-0251:phyxminiproblems-problemphyxmini0251-displacementtimegraph}
  \lean{PhyXMiniProblems.ProblemPhyXMini0251.DisplacementTimeGraph}
  \uses{def:physics:phyx-mini-0251:phyxminiproblems-problemphyxmini0251-lengthquantity, def:physics:phyx-mini-0251:phyxminiproblems-problemphyxmini0251-graphaxis, def:physics:phyx-mini-0251:phyxminiproblems-problemphyxmini0251-axisquantity}
  Labels, scale, and rendering data read from the primary figure. The graph does not print a horizontal-axis unit. timeCoordinateUnit is an arbitrary coherent unit chosen only to take scalar readouts in the governing law; horizontalAxisShowsUnit = false records what the image actually shows
\end{definition}
\begin{proof}
  This is the declared model definition of Displacement Time Graph.
\end{proof}
\begin{definition}[Harmonic Oscillator Graph Setup]
  \label{def:physics:phyx-mini-0251:phyxminiproblems-problemphyxmini0251-harmonicoscillatorgraphsetup}
  \lean{PhyXMiniProblems.ProblemPhyXMini0251.HarmonicOscillatorGraphSetup}
  \uses{def:physics:phyx-mini-0251:phyxminiproblems-problemphyxmini0251-lengthquantity, def:physics:phyx-mini-0251:phyxminiproblems-problemphyxmini0251-timequantity, def:physics:phyx-mini-0251:phyxminiproblems-problemphyxmini0251-angularfrequencyquantity, def:physics:phyx-mini-0251:phyxminiproblems-problemphyxmini0251-velocityquantity, def:physics:phyx-mini-0251:phyxminiproblems-problemphyxmini0251-displacementtimegraph}
  The oscillator quantities referred to by the problem and its graph. position and velocity are independent dimensionful observables. They are connected to amplitudeXm, angularFrequency, and phaseConstantRadians only by the governing-law premise below; no requested phase value is stored in this setup
\end{definition}
\begin{proof}
  This is the declared model definition of Harmonic Oscillator Graph Setup.
\end{proof}
\begin{definition}[Matches Oscillator Figure]
  \label{def:physics:phyx-mini-0251:phyxminiproblems-problemphyxmini0251-matchesoscillatorfigure}
  \lean{PhyXMiniProblems.ProblemPhyXMini0251.MatchesOscillatorFigure}
  \uses{def:physics:phyx-mini-0251:phyxminiproblems-problemphyxmini0251-lengthreadout, def:physics:phyx-mini-0251:phyxminiproblems-problemphyxmini0251-timereadout, def:physics:phyx-mini-0251:phyxminiproblems-problemphyxmini0251-velocityreadout, def:physics:phyx-mini-0251:phyxminiproblems-problemphyxmini0251-harmonicoscillatorgraphsetup}
  Primary-image evidence. The extrema touch the marked levels x\_s and -x\_s, so the amplitude readout is x\_s. The curve crosses the vertical axis one of the three grid intervals below equilibrium and is descending. These are calibrated figure readouts, not assumptions about the requested phase constant
\end{definition}
\begin{proof}
  This is the declared model definition of Matches Oscillator Figure.
\end{proof}
\begin{definition}[Has Physical Harmonic Parameters]
  \label{def:physics:phyx-mini-0251:phyxminiproblems-problemphyxmini0251-hasphysicalharmonicparameters}
  \lean{PhyXMiniProblems.ProblemPhyXMini0251.HasPhysicalHarmonicParameters}
  \uses{def:physics:phyx-mini-0251:phyxminiproblems-problemphyxmini0251-lengthreadout, def:physics:phyx-mini-0251:phyxminiproblems-problemphyxmini0251-angularfrequencyreadout, def:physics:phyx-mini-0251:phyxminiproblems-problemphyxmini0251-harmonicoscillatorgraphsetup}
  Nondegeneracy and the conventional positive phase branch. The full-turn bound is a phase convention; it does not fix the quadrant or numerical answer. The figure's negative initial velocity selects the quadrant through the velocity law below
\end{definition}
\begin{proof}
  This is the declared model definition of Has Physical Harmonic Parameters.
\end{proof}
\begin{definition}[Satisfies Cosine Harmonic Motion]
  \label{def:physics:phyx-mini-0251:phyxminiproblems-problemphyxmini0251-satisfiescosineharmonicmotion}
  \lean{PhyXMiniProblems.ProblemPhyXMini0251.SatisfiesCosineHarmonicMotion}
  \uses{def:physics:phyx-mini-0251:phyxminiproblems-problemphyxmini0251-timequantity, def:physics:phyx-mini-0251:phyxminiproblems-problemphyxmini0251-lengthreadout, def:physics:phyx-mini-0251:phyxminiproblems-problemphyxmini0251-timereadout, def:physics:phyx-mini-0251:phyxminiproblems-problemphyxmini0251-angularfrequencyreadout, def:physics:phyx-mini-0251:phyxminiproblems-problemphyxmini0251-velocityreadout, def:physics:phyx-mini-0251:phyxminiproblems-problemphyxmini0251-harmonicoscillatorgraphsetup}
  The governing simple-harmonic-motion laws in the sign convention printed in the question. The velocity equation is the physical time derivative of the position equation. Neither field assumes the requested phase or an answer choice. Physlib's HarmonicOscillator.AmplitudePhase uses the alternative convention A cos (omega t - phi) and scalar coordinates. This local interface keeps the problem's plus-sign convention and connects dimensionful observables via coherent unit readouts
\end{definition}
\begin{proof}
  This is the declared model definition of Satisfies Cosine Harmonic Motion.
\end{proof}
\begin{definition}[Answer Choice]
  \label{def:physics:phyx-mini-0251:phyxminiproblems-problemphyxmini0251-answerchoice}
  \lean{PhyXMiniProblems.ProblemPhyXMini0251.AnswerChoice}
  Labels of the four answer choices printed with the problem
\end{definition}
\begin{proof}
  This is the declared model definition of Answer Choice.
\end{proof}
\begin{definition}[answer Choice Radians]
  \label{def:physics:phyx-mini-0251:phyxminiproblems-problemphyxmini0251-answerchoiceradians}
  \lean{PhyXMiniProblems.ProblemPhyXMini0251.answerChoiceRadians}
  \uses{def:physics:phyx-mini-0251:phyxminiproblems-problemphyxmini0251-answerchoice}
  Phase value in radians printed beside each answer choice
\end{definition}
\begin{proof}
  This is the declared model definition of answer Choice Radians.
\end{proof}
\begin{definition}[recorded Answer Choice]
  \label{def:physics:phyx-mini-0251:phyxminiproblems-problemphyxmini0251-recordedanswerchoice}
  \lean{PhyXMiniProblems.ProblemPhyXMini0251.recordedAnswerChoice}
  \uses{def:physics:phyx-mini-0251:phyxminiproblems-problemphyxmini0251-answerchoice}
  The answer label recorded in the supplied dataset metadata
\end{definition}
\begin{proof}
  This is the declared model definition of recorded Answer Choice.
\end{proof}
\begin{definition}[Matches Answer To Nearest Hundredth]
  \label{def:physics:phyx-mini-0251:phyxminiproblems-problemphyxmini0251-matchesanswertonearesthundredth}
  \lean{PhyXMiniProblems.ProblemPhyXMini0251.MatchesAnswerToNearestHundredth}
  \uses{def:physics:phyx-mini-0251:phyxminiproblems-problemphyxmini0251-answerchoice, def:physics:phyx-mini-0251:phyxminiproblems-problemphyxmini0251-answerchoiceradians}
  Agreement with a displayed phase to the nearest hundredth of a radian
\end{definition}
\begin{proof}
  This is the declared model definition of Matches Answer To Nearest Hundredth.
\end{proof}
\begin{lemma}[initial phase cosine eq neg one third]
  \label{lem:physics:phyx-mini-0251:phyxminiproblems-problemphyxmini0251-initial-phase-cosine-eq-neg-one-third}
  \lean{PhyXMiniProblems.ProblemPhyXMini0251.initial_phase_cosine_eq_neg_one_third}
  \uses{def:physics:phyx-mini-0251:phyxminiproblems-problemphyxmini0251-harmonicoscillatorgraphsetup, def:physics:phyx-mini-0251:phyxminiproblems-problemphyxmini0251-matchesoscillatorfigure, def:physics:phyx-mini-0251:phyxminiproblems-problemphyxmini0251-hasphysicalharmonicparameters, def:physics:phyx-mini-0251:phyxminiproblems-problemphyxmini0251-satisfiescosineharmonicmotion}
  The initial graph ordinate and the cosine law give cos phi = -1/3
\end{lemma}
\begin{proof}
  The conclusion follows from the governing relations and figure readouts named in the statement of initial phase cosine eq neg one third.
\end{proof}
\begin{lemma}[phase lies in second quadrant]
  \label{lem:physics:phyx-mini-0251:phyxminiproblems-problemphyxmini0251-phase-lies-in-second-quadrant}
  \lean{PhyXMiniProblems.ProblemPhyXMini0251.phase_lies_in_second_quadrant}
  \uses{def:physics:phyx-mini-0251:phyxminiproblems-problemphyxmini0251-harmonicoscillatorgraphsetup, def:physics:phyx-mini-0251:phyxminiproblems-problemphyxmini0251-matchesoscillatorfigure, def:physics:phyx-mini-0251:phyxminiproblems-problemphyxmini0251-hasphysicalharmonicparameters, def:physics:phyx-mini-0251:phyxminiproblems-problemphyxmini0251-satisfiescosineharmonicmotion}
  The descending curve gives sin phi > 0; together with negative cosine and the positive full-turn convention, this places the phase in quadrant II
\end{lemma}
\begin{proof}
  The conclusion follows from the governing relations and figure readouts named in the statement of phase lies in second quadrant.
\end{proof}
\begin{lemma}[phase constant eq arccos neg one third]
  \label{lem:physics:phyx-mini-0251:phyxminiproblems-problemphyxmini0251-phase-constant-eq-arccos-neg-one-third}
  \lean{PhyXMiniProblems.ProblemPhyXMini0251.phase_constant_eq_arccos_neg_one_third}
  \uses{def:physics:phyx-mini-0251:phyxminiproblems-problemphyxmini0251-harmonicoscillatorgraphsetup, def:physics:phyx-mini-0251:phyxminiproblems-problemphyxmini0251-matchesoscillatorfigure, def:physics:phyx-mini-0251:phyxminiproblems-problemphyxmini0251-hasphysicalharmonicparameters, def:physics:phyx-mini-0251:phyxminiproblems-problemphyxmini0251-satisfiescosineharmonicmotion}
  On the selected quadrant, the exact phase is the principal inverse cosine
\end{lemma}
\begin{proof}
  The conclusion follows from the governing relations and figure readouts named in the statement of phase constant eq arccos neg one third.
\end{proof}
\begin{lemma}[arccos neg one third matches choice C]
  \label{lem:physics:phyx-mini-0251:phyxminiproblems-problemphyxmini0251-arccos-neg-one-third-matches-choice-c}
  \lean{PhyXMiniProblems.ProblemPhyXMini0251.arccos_neg_one_third_matches_choice_C}
  \uses{def:physics:phyx-mini-0251:phyxminiproblems-problemphyxmini0251-matchesanswertonearesthundredth}
  The exact inverse-cosine value rounds to 1.91 rad, displayed as choice C
\end{lemma}
\begin{proof}
  The conclusion follows from the governing relations and figure readouts named in the statement of arccos neg one third matches choice C.
\end{proof}
% --- Archon declaration topology end ---
% --- Archon physics formalization source end ---
